\section{National-Scale Findings}

We applied recursive bisection to all 50 states across three census decades (2000, 2010, 2020), generating 1,305 congressional districts total (435 per decade). Six findings emerge from this national-scale demonstration.

\subsection{Regional Context}

Algorithmic redistricting performance varies substantially by U.S. Census region, reflecting differences in demographic composition, geographic clustering, and enacted plan characteristics. Table~\ref{tab:regional-summary} presents key metrics by region (Northeast, Midwest, South, West), establishing geographic context for the findings that follow.

% Regional Summary Table: Key Metrics by U.S. Census Region
% Data aggregated from Papers 1-11 (state-level results rolled up to regions)
% Northeast: CT, ME, MA, NH, RI, VT, NJ, NY, PA (9 states, 75 districts)
% Midwest: IL, IN, MI, OH, WI, IA, KS, MN, MO, NE, ND, SD (12 states, 107 districts)
% South: DE, DC, FL, GA, MD, NC, SC, VA, WV, AL, KY, MS, TN, AR, LA, OK, TX (17 entities, 179 districts)
% West: AZ, CO, ID, MT, NV, NM, UT, WY, AK, CA, HI, OR, WA (13 states, 74 districts)

\begin{table}[htbp]
\centering
\caption{Regional Summary: Algorithmic Redistricting Performance by U.S. Census Region}
\label{tab:regional-summary}
\small
\begin{tabular}{lcccc}
\toprule
\textbf{Metric} & \textbf{Northeast} & \textbf{Midwest} & \textbf{South} & \textbf{West} \\
\midrule
\textit{Basic Characteristics} & & & & \\
States & 9 & 12 & 17 & 13 \\
Congressional Districts & 75 & 107 & 179 & 74 \\
Mean Population Density (per km²) & 342 & 89 & 78 & 42 \\
Mean Non-White (\%) & 28 & 24 & 38 & 35 \\
\midrule
\textit{Compactness (Polsby-Popper, mean)} & & & & \\
Algorithmic Plans & 0.38 & 0.33 & 0.29 & 0.31 \\
Enacted Plans & 0.24 & 0.21 & 0.19 & 0.22 \\
Improvement (\%) & +58 & +57 & +53 & +41 \\
\midrule
\textit{VRA Compliance (MM Districts, total)} & & & & \\
Algorithmic Plans & 8 & 18 & 84 & 27 \\
Enacted Plans & 4 & 9 & 45 & 10 \\
Surplus & +4 & +9 & +39 & +17 \\
\midrule
\textit{42\% Threshold (states meeting)} & & & & \\
Algorithmic Plans & 0 / 9 & 1 / 12 & 9 / 17 & 3 / 13 \\
Threshold Range & 18-32\% & 15-38\% & 24-61\% & 22-51\% \\
\midrule
\textit{Efficiency Gap (\%, mean 2016-2020)} & & & & \\
Algorithmic Plans & $-3.5$ & $-2.9$ & $-3.2$ & $-3.3$ \\
Enacted Plans & $+4.2$ & $+6.8$ & $+5.6$ & $+5.1$ \\
Difference (pp) & 7.7 & 9.7 & 8.8 & 8.4 \\
\midrule
\textit{Temporal Stability (2010→2020, \% tracts retained)} & & & & \\
Recursive Bisection & 82 & 79 & 78 & 81 \\
N-way Partitioning & 68 & 65 & 64 & 67 \\
Advantage (pp) & +14 & +14 & +14 & +14 \\
\bottomrule
\end{tabular}
\vspace{0.5em}

\footnotesize
\textit{Notes}: Regional classifications follow U.S. Census Bureau definitions. MM (majority-minority) districts defined as $\geq 50\%$ non-white population. 42\% threshold indicates states where total minority population $\geq 42\%$ (empirical threshold for achieving near-proportional MM representation). Efficiency gap: positive values favor Republicans, negative favor Democrats. Polsby-Popper compactness ranges 0-1 (higher = more compact). Temporal stability measures census tract retention when redistricting across decades. Data aggregated from Papers 1-11; see individual papers for state-level details. South includes DC. Improvement calculated as (Algorithmic - Enacted) / Enacted × 100.

\end{table}


\textbf{Northeast} (9 states, 75 districts): High population density and urban concentration produce strong compactness scores (mean Polsby-Popper: 0.38) and modest VRA opportunities (8 majority-minority districts, primarily in New York and New Jersey). Geographic Democratic clustering yields algorithmic efficiency gaps of $-3.5\%$ despite low enacted bias ($+4.2\%$), reflecting urban packing inherent to compact districting.

\textbf{Midwest} (12 states, 107 districts): Mixed urban-rural geography produces moderate compactness (0.33) and significant VRA opportunities (18 MM districts, concentrated in Illinois, Michigan, Ohio). This region exhibits the \emph{largest} enacted partisan bias: efficiency gaps of $+6.8\%$ enacted versus $-2.9\%$ algorithmic (9.7 pp difference), driven by aggressive gerrymandering in Rust Belt states following 2010 redistricting.

\textbf{South} (16 states + DC, 179 districts): High minority populations (mean 38\% non-white) enable substantial VRA compliance (84 MM districts, 61\% of national total), with multiple states exceeding the 42\% demographic threshold. Enacted Republican advantage ($+5.6\%$ efficiency gap) substantially exceeds algorithmic Democratic advantage ($-3.2\%$), yielding 8.8 pp improvement. Lower compactness scores (0.29) reflect irregular state boundaries and dispersed rural geography.

\textbf{West} (13 states, 74 districts): Geographic diversity---from Alaska's vast expanse to California's dense metros---produces highly variable compactness (0.31 mean, 0.42 standard deviation). Moderate VRA performance (27 MM districts) reflects Hispanic concentration in Southwest states. Enacted bias ($+5.1\%$) moderately exceeds algorithmic bias ($-3.3\%$), with 8.4 pp improvement.

These regional patterns demonstrate that algorithmic redistricting consistently reduces partisan bias across all regions, though the \emph{magnitude} of improvement varies: largest in Midwest (9.7 pp), smallest in Northeast (7.7 pp). VRA compliance concentrates in South (61\% of MM districts) despite the region comprising 41\% of total districts, reflecting demographics rather than algorithmic design. Compactness inversely correlates with state area: dense Northeastern states achieve higher scores than sprawling Western states, independent of algorithmic choices~\cite{deluca2026compactness}.

\subsection{Finding 1: Technical Feasibility at National Scale}

Algorithmic redistricting meets constitutional and practical requirements. Every district achieved geographic contiguity---the algorithm guarantees this by construction through graph connectivity. Compactness scores (Polsby-Popper metric) improved 56\% over unweighted baseline partitions, demonstrating that edge-weighting successfully balances population equality with geometric optimization~\cite{deluca2026edgeweighted}.

Population equality constraints are user-configurable. The Supreme Court requires ``absolute equality'' for congressional districts absent extraordinary justification (\emph{Karcher v. Daggett}, 1983), with courts scrutinizing deviations exceeding ±0.5\%. The METIS algorithm accepts population balance parameters enabling ±0.5\% or tighter tolerances. For policy adoption, we recommend ±0.5\% maximum deviation for congressional redistricting to ensure compliance with \emph{Karcher}'s stringent standard. State legislative districts permit greater flexibility (~10\% total deviation per \emph{Brown v. Thomson}, 1983)~\cite{deluca2026recursive}.

Computational efficiency enables policy adoption: individual states partition in seconds, and the full 50-state redistricting completes in under two hours on standard hardware. This speed permits ensemble generation---producing thousands of alternative plans to assess statistical properties---impossible with human mapmaking. The method scales from Vermont's single district to California's 52 districts without algorithmic modification.

\subsection{Finding 2: VRA Compliance Exceeds Enacted Plans}

Algorithmic plans created 137 majority-minority districts compared to 68 in enacted maps, a surplus of 69 districts (+101\% increase). This demonstrates that neutral algorithmic methods \emph{can} produce substantial minority representation through geographic optimization alone, without explicit racial targeting. Importantly, this analysis addresses only the \emph{first prong} of the \emph{Thornburg v. Gingles} (1986) test---whether minority groups are sufficiently large and geographically compact to constitute district majorities. Full Section 2 compliance requires two additional showings: that the minority group is politically cohesive, and that white bloc voting usually defeats minority-preferred candidates (prongs 2 and 3)~\cite{deluca2026vra}.

Thus, these 137 districts represent \emph{opportunity districts}---locations where minority populations exceed thresholds necessary for electoral success---rather than \emph{performing districts} confirmed through polarized voting analysis. The finding's significance lies in showing that geographic compactness constraints permit more minority-majority districts than currently enacted, not in claiming that all 137 districts are legally required under Section 2. State-level examples illustrate: Alabama, with zero majority-minority districts in its enacted 2020 plan (later ruled unconstitutional), generates two algorithmically; Georgia increases from five to six; Louisiana from one to two~\cite{deluca2026vra}.

Regional patterns reveal geographic concentration of VRA opportunities. The \textbf{South} accounts for 84 of 137 algorithmic MM districts (61\%), despite comprising 41\% of total congressional seats, reflecting high minority populations (mean 38\% non-white) and favorable spatial clustering. Key Southern states include Texas (14 MM districts algorithmically vs. 8 enacted), Florida (11 vs. 5), and Georgia (6 vs. 5). The \textbf{Midwest} contributes 18 MM districts (13\%), concentrated in Illinois, Michigan, and Ohio. The \textbf{West} produces 27 districts (20\%), primarily in California (15 MM districts) and Arizona (3). The \textbf{Northeast} yields 8 districts (6\%), concentrated in New York (5) and New Jersey (2). This geographic distribution reflects demography rather than algorithmic bias: regions with larger minority populations and greater residential clustering naturally enable more MM districts under neutral redistricting principles~\cite{deluca2026vra}.

\subsection{Finding 3: The 42\% Demographic Threshold}

Empirical analysis reveals a feasibility threshold: states with minority populations exceeding 42\% achieve near-proportional minority representation (within 5 percentage points), while states below 37\% minority fail regardless of algorithmic sophistication. This 42\% threshold provides empirical benchmarks for the \emph{geographic compactness component} of the first \emph{Gingles} prong---whether minority populations are sufficiently large and geographically clustered to permit majority-minority districts under neutral redistricting principles~\cite{deluca2026threshold}.

The threshold reflects geographic clustering, not population percentage alone. Mississippi (38\% Black) creates two majority-Black districts through favorable spatial distribution, while New York (24\% Hispanic) creates only three Hispanic-majority districts despite larger absolute population. Courts adjudicating Section 2 VRA claims currently assess geographic compactness through expert testimony and illustrative maps; the 42\% threshold offers a complementary quantitative tool derived from national-scale analysis. However, crossing this threshold indicates geometric \emph{possibility}, not legal \emph{obligation}---full Section 2 liability requires demonstrating vote dilution through all three \emph{Gingles} prongs~\cite{deluca2026threshold}.

Regional distribution of states meeting the 42\% threshold reveals stark geographic patterns. The \textbf{South} dominates with 9 of 13 qualifying states (69\%): Alabama, Georgia, Louisiana, Maryland, Mississippi, North Carolina, South Carolina, Texas, and Virginia. These states range from 42\% to 61\% minority population, with Mississippi (61\%), Georgia (52\%), and Maryland (51\%) exhibiting the highest concentrations. The \textbf{West} contributes 3 states: California (51\%), New Mexico (50\%), and Nevada (45\%). The \textbf{Midwest} provides 1 state: Illinois (43\%). The \textbf{Northeast} contributes 0 states, with New York's 38\% minority population falling just below threshold. This geographic concentration reflects historical settlement patterns, the Great Migration, and recent Hispanic immigration to Southwest states. Notably, all 9 Southern states meeting the threshold achieve near-proportional MM representation algorithmically (within 5 percentage points of minority population share), confirming the threshold's predictive validity~\cite{deluca2026threshold}.

\subsection{Finding 4: Partisan Patterns Reflect Geography}

Without accessing partisan data, the algorithm produces districts with 56.5\% Democratic vote share (estimated from precinct-level election results post-facto). This partisan asymmetry arises from geographic sorting: Democratic voters concentrate in urban cores while Republican voters disperse across suburbs and rural areas. Recursive bisection, optimizing for geographic compactness, unavoidably packs urban Democrats while splitting rural Republicans---the opposite pattern emerges from geography, not algorithmic intent~\cite{chen2013unintentional,deluca2026edgeweighted}.

This finding strengthens the impossibility defense. Critics might argue that ``neutral'' algorithms still produce partisan advantage, but the advantage reflects voter distribution, not gerrymandering. Human mapmakers can amplify or reverse geographic patterns by strategic cracking and packing; algorithms cannot. The efficiency gap---a metric measuring ``wasted votes''---measures geographic artifacts when the mapmaker cannot see partisan data~\cite{deluca2026edgeweighted}.

\subsection{Finding 5: Reduced Partisan Bias Compared to Enacted Plans}

Quantitative partisan fairness analysis reveals that algorithmic redistricting substantially reduces partisan bias relative to enacted maps. Using the \emph{efficiency gap}---which measures wasted votes (ballots for losing candidates plus surplus votes beyond the threshold needed to win)---we compared algorithmic and enacted plans across competitive states for 2016-2020 elections. Algorithmic plans exhibit a mean efficiency gap of $-3.2\%$ (slight Democratic advantage), while enacted plans show $+5.1\%$ (Republican advantage), a difference of 8.3 percentage points. This represents a 62\% reduction in partisan bias under algorithmic redistricting~\cite{stephanopoulos2015partisan,deluca2026edgeweighted}.

The negative efficiency gap in algorithmic plans reflects geographic sorting: urban Democratic concentration creates unavoidable packing when optimizing for compactness. However, this geographic bias ($-3.2\%$) is substantially smaller than the bias in enacted plans ($+5.1\%$), suggesting that human mapmakers amplify rather than merely reflect geographic asymmetries. Importantly, algorithmic plans approach but do not achieve perfect proportionality---the $-3.2\%$ efficiency gap indicates Democrats win 2-3 more seats nationally than strict proportionality would predict. This modest asymmetry represents the minimum partisan effects achievable under geographic constraints with compact districts~\cite{chen2013unintentional}.

Proportionality analysis confirms these patterns. Across competitive states where Democrats averaged 51.3\% of two-party vote share (2016-2020), algorithmic plans produce 54.1\% Democratic seats (proportionality gap: $+2.8$ pp) while enacted plans produce 46.8\% Democratic seats (proportionality gap: $-4.5$ pp). Mean-median difference analysis shows algorithmic plans pack Democrats modestly ($+0.8$ pp mean-median gap) while enacted plans pack Democrats severely ($+4.1$ pp). Seats-votes curve analysis reveals algorithmic plans produce near-symmetric responsiveness (elasticity: 2.8) while enacted plans show dampened responsiveness favoring incumbents (elasticity: 2.1). These multiple metrics converge on the conclusion that algorithmic redistricting substantially improves proportionality compared to enacted plans, though geographic constraints prevent perfect proportionality~\cite{deluca2026efficiency}.

State-level variation reveals regional patterns. Rust Belt states (Pennsylvania, Wisconsin, Michigan) show the largest differences: enacted plans average $+7.2\%$ efficiency gaps (strong Republican advantage) while algorithmic plans average $-2.8\%$ (modest Democratic advantage). Sunbelt states (Arizona, Georgia, North Carolina) show smaller but consistent patterns: enacted plans favor Republicans by $+4.3\%$, algorithmic plans by $-3.1\%$. These findings demonstrate that algorithmic redistricting reduces but cannot eliminate partisan asymmetries driven by residential geography~\cite{rodden2019cities,deluca2026edgeweighted}.

\subsection{Finding 6: Temporal Stability Across Census Cycles}

Recursive bisection maintains 80\% census tract retention when redistricting from 2010 to 2020 boundaries, a 14-percentage-point advantage over n-way partitioning methods. The algorithm's hierarchical structure---districts emerge from nested binary divisions---creates geographic scaffolding that persists across demographic change. Statistical analysis reveals variance between geographic regions exceeds temporal variance by 3.2×, indicating district boundaries reflect durable geographic structure rather than decade-specific population quirks~\cite{deluca2026temporal}.

Regional patterns show consistent stability advantage across all four Census regions. The \textbf{Northeast} achieves 82\% tract retention (vs. 68\% for n-way, +14 pp), the \textbf{Midwest} 79\% (vs. 65\%, +14 pp), the \textbf{South} 78\% (vs. 64\%, +14 pp), and the \textbf{West} 81\% (vs. 67\%, +14 pp). The uniform 14 percentage point advantage demonstrates that hierarchical bisection's stability benefit derives from algorithmic structure rather than regional geography. Slight variation in absolute retention rates (78-82\%) reflects differences in demographic volatility: Sunbelt states experiencing rapid population growth show marginally lower retention than stable Rust Belt states, but the recursive advantage persists regardless~\cite{deluca2026temporal}.

Stability matters for democratic accountability: voters can build relationships with representatives, and representatives can develop district expertise when boundaries change minimally. Temporal consistency validates the method as policy infrastructure rather than one-time experiment, suitable for repeated application across multiple census cycles~\cite{deluca2026temporal}.
